\documentclass[UTF8,zihao=-4,a4paper]{ctexart}

\usepackage{geometry}
\usepackage{fontspec}
\usepackage{xeCJK}
\usepackage{xcolor}
\usepackage{graphicx}
\usepackage{booktabs}
\usepackage{tabularx}
\usepackage{array}
\usepackage{longtable}
\usepackage{amsmath}
\usepackage{tikz}
\usepackage{enumitem}
\usepackage{listings}
\usepackage{caption}
\usepackage{fancyhdr}
\usepackage{titlesec}
\usepackage{etoolbox}
\usepackage[hidelinks]{hyperref}

\usetikzlibrary{arrows.meta,positioning,fit,calc,shapes.geometric}

\geometry{left=2.65cm,right=2.65cm,top=2.55cm,bottom=2.65cm}
\xeCJKsetup{CJKecglue={\hskip 0.10em plus 0.03em minus 0.02em}}
\IfFontExistsTF{Times New Roman}{\setmainfont{Times New Roman}}{\setmainfont{TeX Gyre Termes}}
\IfFontExistsTF{Arial}{\setsansfont{Arial}}{\setsansfont{TeX Gyre Heros}}
\IfFontExistsTF{Consolas}{\setmonofont{Consolas}}{\setmonofont{Courier New}}
\IfFontExistsTF{SimSun}{\setCJKmainfont{SimSun}[BoldFont={SimHei}]}{\setCJKmainfont{FandolSong-Regular}}
\IfFontExistsTF{Microsoft YaHei}{\setCJKsansfont{Microsoft YaHei}[BoldFont={Microsoft YaHei}]}{\setCJKsansfont{FandolHei-Regular}}
\IfFontExistsTF{FangSong}{\setCJKmonofont{FangSong}}{\setCJKmonofont{FandolFang-Regular}}

\definecolor{moonblue}{HTML}{0B5E7A}
\definecolor{mooncyan}{HTML}{26A6B8}
\definecolor{moongray}{HTML}{F3F6F8}
\definecolor{moonorange}{HTML}{E07A35}
\definecolor{moongreen}{HTML}{2D8A58}
\definecolor{moonred}{HTML}{B5423A}
\definecolor{codegray}{HTML}{F6F8FA}

\hypersetup{
  pdftitle={CCF2026 Mooncake Store碎片感知分配优化技术报告},
  pdfauthor={Lorry1024},
  pdfsubject={Mooncake Store KVCache存储设计和性能优化},
  pdfkeywords={Mooncake Store,KVCache,fragmentation aware,allocation strategy}
}

\captionsetup{font={small,sf},labelfont={bf,sf},labelsep=quad}
\setlength{\parindent}{2em}
\setlength{\parskip}{0.25em}
\setlength{\headheight}{14pt}
\emergencystretch=3em
\hfuzz=1pt
\renewcommand{\arraystretch}{1.34}
\setlength{\tabcolsep}{5pt}
\setlength{\arrayrulewidth}{0.45pt}
\setlength{\abovedisplayskip}{0.72em}
\setlength{\belowdisplayskip}{0.72em}
\setlength{\abovedisplayshortskip}{0.55em}
\setlength{\belowdisplayshortskip}{0.55em}
\makeatletter
\setlength{\@fptop}{0pt}
\makeatother
\BeforeBeginEnvironment{tabular}{\begingroup\small\sffamily}
\AfterEndEnvironment{tabular}{\endgroup}
\BeforeBeginEnvironment{tabularx}{\begingroup\small\sffamily}
\AfterEndEnvironment{tabularx}{\endgroup}
\BeforeBeginEnvironment{longtable}{\begingroup\small\sffamily}
\AfterEndEnvironment{longtable}{\endgroup}
\AtBeginEnvironment{align}{\small}
\setlist[itemize]{leftmargin=2.2em,itemsep=0.2em,topsep=0.3em}
\setlist[enumerate]{leftmargin=2.2em,itemsep=0.2em,topsep=0.3em}

\titleformat{\section}{\Large\bfseries\sffamily\color{black}}{\thesection}{0.8em}{}
\titleformat{\subsection}{\large\bfseries\sffamily\color{black}}{\thesubsection}{0.7em}{}
\titleformat{\subsubsection}{\normalsize\bfseries\sffamily\color{black}}{\thesubsubsection}{0.6em}{}

\pagestyle{fancy}
\fancyhf{}
\fancyhead[C]{\small Mooncake Store碎片感知分配优化技术报告}
\fancyfoot[C]{\small\thepage}
\renewcommand{\headrulewidth}{0.45pt}
\renewcommand{\footrulewidth}{0.45pt}

\lstdefinestyle{codestyle}{
  basicstyle=\ttfamily\small,
  backgroundcolor=\color{codegray},
  frame=single,
  rulecolor=\color{moonblue!35},
  breaklines=true,
  columns=fullflexible,
  keepspaces=true,
  showstringspaces=false,
  tabsize=2
}
\lstset{style=codestyle}

\newcommand{\artifact}[1]{\texttt{#1}}
\newcommand{\smallartifact}[1]{{\ttfamily\footnotesize #1}}
\newcommand{\strategy}{\texttt{fragmentation\_aware}}
\newcommand{\baseline}{\texttt{free\_ratio\_first}}
\newcommand{\mvar}[1]{\mathit{#1}}
\newcommand{\thfirst}[1]{\multicolumn{1}{|c|}{\bfseries #1}}
\newcommand{\thcell}[1]{\multicolumn{1}{c|}{\bfseries #1}}
\newcommand{\splitpath}[2]{\begin{tabular}[t]{@{}l@{}}\ttfamily\scriptsize #1\\\ttfamily\scriptsize #2\end{tabular}}
\newcommand{\splitpaththree}[3]{\begin{tabular}[t]{@{}l@{}}\ttfamily\scriptsize #1\\\ttfamily\scriptsize #2\\\ttfamily\scriptsize #3\end{tabular}}
\newcolumntype{Y}{>{\raggedright\arraybackslash}X}
\newcolumntype{L}[1]{>{\raggedright\arraybackslash}p{#1}}
\newcolumntype{C}[1]{>{\centering\arraybackslash}p{#1}}
\newcolumntype{B}[1]{>{\centering\arraybackslash\bfseries}p{#1}}

\begin{document}

\begin{titlepage}
  \centering
  \vspace*{1.1cm}
  {\sffamily\bfseries\zihao{2}\color{black}CCF2026开源创新大赛\par}
  \vspace{0.35cm}
  {\sffamily\bfseries\zihao{3}\color{black}Mooncake赛题2初赛技术报告\par}
  \vspace{2.0cm}
  {\bfseries\zihao{1}Mooncake Store碎片感知分配优化\par}
  \vspace{0.55cm}
  {\Large Fragmentation-Aware Allocation for Mooncake Store\par}
  \vspace{2.2cm}
  \begin{tabular}{|>{\sffamily\bfseries\centering\arraybackslash}p{3.2cm}|>{\raggedright\arraybackslash}p{9.3cm}|}
    \hline
    参赛方向 & 赛题2：优化Mooncake Store吞吐性能、高可用功能和可扩展性，优化SGLang HiCache+Mooncake Store性能 \\ \hline
    作品名称 & Mooncake Store碎片感知分配策略 \\ \hline
    参赛仓库 & \url{https://github.com/Lorry1024/ccf2026-mooncake-fragmentation-aware} \\ \hline
    上游PR & \url{https://github.com/kvcache-ai/Mooncake/pull/2797} \\ \hline
    PR头提交 & \artifact{0123fa1 Fix fragmentation-aware allocation test setup} \\ \hline
    报告日期 & 2026年7月9日 \\ \hline
  \end{tabular}
  \vfill
  {\small 本报告与仓库中的patch、复现实验、验证日志和展示PPT共同构成初赛提交材料。\par}
\end{titlepage}

\pagenumbering{Roman}
\section*{摘要}

Mooncake Store是Mooncake体系中支撑KVCache对象存储和复用的关键组件。现有\baseline{}策略根据segment总空闲比例进行排序，能够改善整体空间利用均衡，但在混合大小KVCache对象和长时间运行的场景下，总空闲空间可能被切分为多个小空洞，无法真实反映当前大对象是否可以被连续放入。本文提出并实现\strategy{}策略，在保持默认行为不变和有界候选采样不变的前提下，将最大连续空闲区域、连续性比例和当前请求可容纳性纳入排序逻辑，使Master优先选择能够直接满足请求的segment。

本作品已形成Mooncake上游draft PR，当前PR头提交\artifact{0123fa1}已通过GitHub Actions，结果为26个检查成功、1个检查跳过。确定性专题仿真显示，在6个大对象请求上，\baseline{}的首选segment可直接容纳次数为0/6，而\strategy{}为6/6；fallback尝试次数从11降至0；平均候选segment数量保持5.00；本地单次决策耗时从121.00ns增加到138.93ns，额外开销约17.93ns。该结果说明，新策略能够以很小的决策成本提升碎片化场景下的分配稳定性。

\noindent\textbf{关键词：}Mooncake Store；KVCache；碎片感知；连续空闲区域；分配策略；SGLang HiCache

\newpage
\section*{模块完成情况}

\begin{table}[h]
  \centering
  \caption{初赛阶段模块完成情况}
  \begin{tabularx}{\textwidth}{|B{3.2cm}|Y|}
    \hline
    \thfirst{模块} & \thcell{完成情况} \\ \hline
    问题分析 & 明确\baseline{}在碎片化大对象分配中的误判模式，形成最小复现和专题仿真。 \\ \hline
    策略设计 & 设计\strategy{}排序规则，优先考虑\artifact{can\_fit}、连续性评分和最大连续空闲区域。 \\ \hline
    代码实现 & 在Mooncake Store中新增策略类型、策略类、配置解析、启动参数和benchmark矩阵。 \\ \hline
    测试验证 & 新增碎片化单元测试，保留preferred segment语义，PR通过上游GitHub Actions。 \\ \hline
    材料整理 & 完成中文Markdown材料、LaTeX技术报告、展示PPT、复现实验源码和提交包。 \\ \hline
    待完成项 & 5分钟展示视频尚未录制；真实RDMA和SGLang HiCache端到端benchmark作为后续工作。 \\ \hline
  \end{tabularx}
\end{table}

\newpage
\tableofcontents
\newpage
\pagenumbering{arabic}

\section{赛题理解与项目定位}

Mooncake赛题2要求围绕Mooncake Store吞吐性能、高可用功能、可扩展性以及SGLang HiCache+Mooncake Store性能开展优化。该要求并不只局限于网络传输或外层推理框架，也包含Store后端内部影响吞吐稳定性和资源利用率的关键路径。对KVCache场景而言，Store中对象的分配位置会影响一次写入是否能顺利落地，也会影响后续fallback、重试、淘汰和副本放置压力。

本项目选择的切入点是Mooncake Store分配策略层。该方向具有三个特点：第一，它位于真实Mooncake Store代码路径中，能够通过PR形式提交给上游；第二，它对SGLang HiCache等前端具有后端相关性，因为这些前端将KVCache对象交给Store时需要依赖Store的放置决策；第三，它可以在不重写协议、不引入重依赖、不改变默认行为的条件下完成可验证优化。

\begin{figure}[h]
  \centering
  \resizebox{0.98\textwidth}{!}{%
  \begin{tikzpicture}[
    node distance=1.0cm and 1.1cm,
    box/.style={draw=moonblue,very thick,rounded corners=3pt,align=center,minimum width=2.75cm,minimum height=0.95cm,fill=moongray},
    hot/.style={draw=moonorange,very thick,rounded corners=3pt,align=center,minimum width=2.9cm,minimum height=0.95cm,fill=moonorange!10},
    arrow/.style={-{Latex[length=3mm]},very thick,draw=moonblue}
  ]
    \node[box] (frontend) {推理前端\\SGLang/HiCache};
    \node[box,right=of frontend] (client) {Mooncake Store\\Client};
    \node[hot,right=of client] (master) {Master\\分配策略};
    \node[box,right=of master] (segments) {候选segments\\有界采样};
    \node[box,below=of segments] (alloc) {OffsetAllocator\\CacheLib};
    \node[box,left=of alloc] (replica) {副本写入\\Store backend};
    \draw[arrow] (frontend) -- (client);
    \draw[arrow] (client) -- (master);
    \draw[arrow] (master) -- (segments);
    \draw[arrow] (segments) -- (alloc);
    \draw[arrow] (alloc) -- (replica);
    \draw[arrow] (replica.west) -- ++(-1.1,0) |- (client.south);
  \end{tikzpicture}
  }
  \caption{本项目在Mooncake Store路径中的位置}
\end{figure}

\subsection{技术路线总览}

本项目的技术路线可以概括为“先把碎片化问题显式建模，再把模型压缩为可在现有分配路径中低成本执行的排序信号”。我们没有从修改网络协议、重写存储后端或引入全局整理线程入手，原因是这些方向虽然可能带来更大范围的能力提升，但会显著增加验证复杂度，并且容易影响Mooncake Store已有部署行为。赛题初赛更需要一个边界清晰、证据充分、能够被上游项目评审的优化点。因此，本项目把关注点收敛到Master侧segment选择过程：当Store面对一个已知大小的KVCache对象时，应该优先选择“当前真的能放下”的segment，而不是只选择“总空闲比例看起来更高”的segment。

从工程路径看，完整工作被拆成六步。第一步，阅读Mooncake Store已有分配策略，确认\artifact{random}、\baseline{}等策略的职责边界。第二步，构造最小碎片化反例，证明总空闲比例和连续可用空间之间存在可观察的不一致。第三步，设计\strategy{}评分字段，把最大连续空闲区域、连续性比例和当前请求大小引入候选排序。第四步，把策略接入Mooncake Store正式配置路径，而不是停留在离线脚本或实验代码中。第五步，补充单元测试、benchmark矩阵和文档说明，保证策略可使用、可回滚、可被评审。第六步，通过上游PR和GitHub Actions完成开源项目层面的基本验证。该路线对应赛题2的三个关键约束：优化点位于Mooncake Store主路径，策略开销受有界候选采样约束控制，并且默认行为、preferred/excluded语义和fallback路径均保持不变。

\section{原有机制与问题分析}

\subsection{Mooncake Store分配链路与关键抽象}

Mooncake Store在写入KVCache对象时，需要为对象的每个副本选择可写入的segment。上层调用并不直接关心底层空闲块的形态，而是把对象大小、副本数、preferred segment和excluded segment等约束交给Master侧分配逻辑。Master再调用具体分配策略，从可用segment集合中选出候选，并通过对应allocator完成空间申请。这个过程看似只是“选择一个segment”，但实际承担了三类语义：

\begin{enumerate}
  \item \textbf{容量语义。} 当前segment是否拥有足够空间，决定对象是否可以一次性落地。
  \item \textbf{副本语义。} 多副本写入需要避免重复使用同一segment，并在候选不足时保留best-effort fallback。
  \item \textbf{调度语义。} preferred segment表示上层希望优先使用的放置位置，excluded segment表示本次请求不能选择的位置。
\end{enumerate}

因此，分配策略不是孤立的性能小优化，而是Store吞吐稳定性、资源利用率和高可用副本放置之间的连接点。本项目的接入点在Master侧策略工厂和配置解析路径；被排序的对象是segment；真正判断空间能否占用的是segment内部allocator；preferred segments和excluded segments则分别代表上层显式优先级和本次请求的禁止集合。如果策略选中了“总空闲很多但无法容纳当前对象”的segment，后续allocator仍会拒绝分配，系统只能继续尝试其他segment或走fallback路径。该失败不会直接制造数据错误，但会增加一次无效尝试，并可能放大为锁竞争、重试、尾延迟和副本放置压力。

\subsection{KVCache对象生命周期与碎片来源}

KVCache对象的存储行为和普通定长块写入不同。推理系统中不同请求的上下文长度、batch形态、prefix共享程度和缓存淘汰时机都可能不同，导致写入Store的对象大小呈现明显差异。一个较短请求可能只产生较小的KVCache对象，一个长上下文请求可能产生更大的对象；多轮对话或prefix复用又会让对象生命周期不完全同步。当这些对象交错写入、交错释放时，同一个segment内部就会逐渐形成不规则的空闲区域。

可以把segment想象成一个连续地址空间。初始状态下，segment中空闲区域通常比较完整，任何大小合理的对象都容易落入。运行一段时间后，不同大小对象被插入其中；再运行一段时间后，部分对象被释放，释放位置不一定相邻。此时虽然空闲字节总量可能很多，但这些字节被分散在多个小区间中。碎片来源主要包括对象大小混合、生命周期不同步、多副本写入、前缀复用命中后的非同步淘汰以及长时间运行中的持续分配释放。对allocator而言，当前请求是否可分配，取决于是否存在一个足够大的连续空闲区间，而不是所有小区间相加是否足够。

这一点正是赛题2中“Store吞吐性能”和“SGLang HiCache+Mooncake Store性能”之间的连接处。HiCache等前端把KVCache对象交给Store，并期望后端尽快完成存放和复用。如果Store在碎片化场景下反复优先选择不可容纳的大对象位置，就会增加后端分配路径的不确定性。即使最终能够通过fallback找到可用位置，前面失败的尝试仍然消耗了决策、锁、状态检查和日志路径资源。

\subsection{现有策略的能力边界}

Mooncake Store已有多种分配策略。默认\artifact{random}策略开销低、行为简单，适合通用场景；\baseline{}策略按照总空闲比例排序，目标是让更空闲的segment优先吸收写入，从而改善整体空间利用均衡；\artifact{ssd\_free\_ratio\_first}、\artifact{cxl}和\artifact{local\_first}等策略则面向特定介质或拓扑。就赛题2而言，\baseline{}是最重要的对比基线，因为它已经使用了空间信息，却仍然没有识别segment内部碎片。

\baseline{}的核心假设可以概括为：总空闲比例越高，当前请求越可能分配成功。该假设在对象大小比较均匀、释放行为比较规则、segment内部空闲区较少被切碎时通常成立。但KVCache Store的运行状态并不总是满足这些条件。长上下文推理、不同batch形态、前缀复用和缓存淘汰会让对象大小差异扩大；对象生命周期不同步又会让释放位置分散。最终，一个segment可能拥有较高总空闲字节数，但这些空间被拆成多个小洞，无法提供一个足够大的连续区域。

\subsection{总空闲比例误判的形成机制}

在allocator需要连续空间的前提下，“总空闲空间”和“最大连续空闲区域”是两个不同指标。前者衡量还有多少空间未被使用，后者衡量当前请求能否被一次性放入。碎片化场景下，两者可能明显背离。图\ref{fig:frag-shape}展示了典型情况：上方segment的绿色空闲块数量多、总量不低，但每个块都较小；下方segment的总空闲量略低，却拥有足够大的连续区域。

\begin{figure}[htbp]
  \centering
  \resizebox{0.98\textwidth}{!}{%
  \begin{tikzpicture}[
    block/.style={draw=white,minimum height=0.55cm},
    used/.style={block,fill=moonblue!75},
    free/.style={block,fill=moongreen!72},
    frag/.style={block,fill=moonorange!70},
    label/.style={font=\small\sffamily}
  ]
    \node[label,anchor=east] at (-0.2,1.2) {fragmented};
    \draw[thick] (0,0.9) rectangle (10,1.5);
    \foreach \x/\w/\c in {0/1.0/frag,1.0/1.0/used,2.0/1.0/frag,3.0/1.0/used,4.0/1.0/frag,5.0/1.0/used,6.0/1.0/frag,7.0/1.0/used,8.0/1.0/frag,9.0/1.0/used} {
      \path[\c] (\x,0.9) rectangle ++(\w,0.6);
    }
    \node[label,anchor=west] at (10.2,1.2) {总空闲高，最大连续块不足};

    \node[label,anchor=east] at (-0.2,0.15) {contiguous};
    \draw[thick] (0,-0.15) rectangle (10,0.45);
    \path[used] (0,-0.15) rectangle (6.2,0.45);
    \path[free] (6.2,-0.15) rectangle (10,0.45);
    \node[label,anchor=west] at (10.2,0.15) {总空闲略低，连续块足够};

    \draw[{Latex[length=2.5mm]}-,thick,moonred] (6.2,-0.45) -- (10,-0.45);
    \node[label,moonred] at (8.1,-0.75) {可容纳10MiB请求};
  \end{tikzpicture}
  }
  \caption{总空闲空间与最大连续空闲区域可能不一致}
  \label{fig:frag-shape}
\end{figure}

最小复现数据如表\ref{tab:minimal-case}所示。\baseline{}会优先选择总空闲空间为16MiB的\artifact{fragmented}，因为它的总空闲比例更高；但这个segment的最大连续空闲区域只有8MiB，无法容纳10MiB请求。相反，\artifact{contiguous}虽然总空闲空间只有12MiB，却可以直接完成本次分配。

\begin{table}[htbp]
  \centering
  \caption{最小碎片化复现场景}
  \label{tab:minimal-case}
  \begin{tabular}{|B{3.0cm}|r|r|r|C{2.4cm}|}
    \hline
    \thfirst{segment} & \thcell{总容量} & \thcell{总空闲} & \thcell{最大连续空闲} & \thcell{10MiB请求} \\ \hline
    \artifact{fragmented} & 32MiB & 16MiB & 8MiB & 不能容纳 \\ \hline
    \artifact{contiguous} & 32MiB & 12MiB & 12MiB & 可以容纳 \\ \hline
  \end{tabular}
\end{table}

\subsection{误判带来的性能影响}

该问题的关键不在于最终一定失败，而在于首选候选质量下降。典型链路是：排序阶段把高总空闲但碎片化严重的segment放在前面，allocator阶段发现最大连续区域小于请求大小，当前副本只能继续尝试其他segment；如果多个副本同时处于碎片化状态，失败尝试会被放大。对推理系统而言，这类失败尝试会影响Store写入路径的稳定性，进而影响KVCache对象的后端落盘、复用和淘汰压力。

\subsection{为什么不能只依赖fallback}

一种直观疑问是：既然\baseline{}选错后仍然可以继续尝试其他segment，为什么还需要新增策略？原因在于fallback解决的是“最终能否找到位置”，而不是“首选路径是否稳定”。在高吞吐推理服务中，Store面对的是连续不断的对象写入请求。一次请求多一次失败尝试看似开销有限，但在高并发、多副本、大对象和长时间运行组合下，会形成累计影响。

首先，fallback通常意味着前一次候选已经做过元数据读取、排除集合判断、allocator尝试和失败处理。其次，多副本请求中每个副本都要避免重复使用同一segment，因此一个副本的失败会改变后续副本的候选空间。再次，fallback路径本质上是补救逻辑，它可以提高可用性，但不应成为碎片化场景下的常态路径。更理想的做法是在排序阶段就把明显不可容纳的候选排到后面，让系统优先尝试成功概率更高的segment。也就是说，fallback解决“失败后如何补救”，而碎片感知排序解决“第一次尝试能否更接近正确位置”。

因此，本项目并不是替代fallback，而是减少不必要fallback。保留fallback是为了处理并发状态变化、候选不足、元数据估计不精确等不可避免情况；新增\strategy{}是为了让正常路径更少走到补救路径。这种关系更符合生产系统设计：主路径尽量准确，兜底路径必须存在，但不应成为常态。

\subsection{本项目抽象出的优化需求}

基于上述分析，本项目没有选择重写Store协议或引入复杂全局整理机制，而是把问题收敛为“在现有候选集合内提高排序质量”。这个抽象有明确边界：它不尝试移动已有对象，不做在线压缩，不改变副本协议，也不扩大候选扫描范围；它只利用allocator已经能提供或可近似得到的元数据，让策略在面对大对象请求时先识别“当前是否真的能放下”。因此，优化需求可以拆成四点：

\begin{enumerate}
  \item 在排序中显式加入最大连续空闲区域，而不是只看总空闲比例。
  \item 在当前请求大小已知的情况下，优先选择\artifact{can\_fit=true}的候选。
  \item 保持原有采样规模和fallback语义，避免策略本身成为新的扩展性瓶颈。
  \item 保持配置可选和默认行为不变，便于灰度、回滚和上游接受。
\end{enumerate}

\section{碎片感知分配策略设计}

\subsection{设计目标与非目标}

\strategy{}的目标是提高碎片化场景下的首选segment可用性。具体来说，它希望在混合大小对象场景中优先选择最大连续空闲区域能够容纳当前请求的segment，沿用\baseline{}的有界候选思想，只改变候选排序信号，不扩大搜索范围；同时保留preferred segment、excluded segment、多副本去重和best-effort fallback语义，并作为显式可选策略接入，不改变\artifact{random}默认行为。它的非目标也很明确：不进行对象迁移、在线压缩或全局碎片整理，不修改Mooncake Store客户端协议、metadata协议或后端存储格式，也不把初赛阶段的策略层验证夸大为真实RDMA集群或SGLang HiCache端到端压测结论。设计时刻意区分目标和非目标，是为了避免优化范围失控。

\subsection{策略输入、输出与不变量}

\strategy{}并不是独立调度器，而是Mooncake Store分配策略体系中的一个可选实现。因此，它的输入和输出必须严格服从已有接口约束。策略输入主要包括当前请求的slice大小、副本数量、allocator manager中的segment视图、preferred segment列表和excluded segment集合。策略输出不是全局放置计划，而是一组已成功申请的buffer句柄；每个成功结果都对应一个实际allocator申请。换言之，评分阶段只决定尝试顺序，真正的空间占用仍以allocator返回结果为准。

本项目在设计中保留三个不变量。第一，任何被excluded的segment都不能被选中。第二，同一次多副本分配中，已经成功使用的segment会进入used集合，后续副本不能重复选择。第三，preferred segment具有显式优先级，只要preferred可用并能完成分配，就不进入后续评分排序。这些不变量保证新策略不会改变调用方已经依赖的放置语义。请求大小\artifact{slice\_length}只用于判断候选是否具备足够连续空间，副本数\artifact{replica\_num}只决定最终需要返回多少个buffer；策略内部不会把评分结果当作空间占用事实。

\subsection{候选元数据与评分字段}

策略对每个采样候选segment计算五类字段。这里的关键是把“容量总量”和“空间形态”拆开：\artifact{free\_ratio}继续描述整体空闲程度，\artifact{largest\_free\_region}和\artifact{contiguity\_ratio}描述空闲区域是否集中，\artifact{can\_fit}则把当前请求大小纳入判断。

\begin{table}[htbp]
  \centering
  \caption{候选segment评分字段}
  \label{tab:score-fields}
  \begin{tabularx}{\textwidth}{|B{3.8cm}|L{4.8cm}|Y|}
    \hline
    \thfirst{字段} & \thcell{计算方式} & \thcell{设计含义} \\ \hline
    \artifact{total\_free} & 所有allocator可用空间之和 & 描述segment剩余空间总量 \\ \hline
    \smallartifact{largest\_free\_region} & allocator报告的最大连续空闲区域最大值 & 描述当前对象是否存在可落入的连续块 \\ \hline
    \artifact{free\_ratio} & \artifact{total\_free / total\_capacity} & 保留原有空间均衡信号 \\ \hline
    \smallartifact{contiguity\_ratio} & \smallartifact{largest\_free\_region / total\_free} & 衡量空闲空间是否集中 \\ \hline
    \artifact{can\_fit} & \smallartifact{largest\_free\_region >= request\_size} & 直接判断当前请求是否可被候选容纳 \\ \hline
  \end{tabularx}
\end{table}

在实现中，若某类allocator无法提供精确的最大连续空闲区域，策略退化为使用可用空闲空间作为估计值。这样做的原因是兼容性优先：新策略不能要求所有allocator一次性改造完成。对能提供精确碎片信息的OffsetAllocator，策略会使用真实\artifact{getLargestFreeRegion()}；对无法提供精确值的场景，行为退化为接近总空闲视角，不会阻断已有部署。

\subsection{排序规则}

评分公式如下：

\begin{align}
  \mvar{free\_ratio} &= \frac{\mvar{total\_free}}{\mvar{total\_capacity}} \\
  \mvar{contiguity\_ratio} &= \frac{\mvar{largest\_free\_region}}{\mvar{total\_free}} \\
  \mvar{score} &= 0.70\times \mvar{contiguity\_ratio}+0.30\times \mvar{free\_ratio}
\end{align}

其中0.70和0.30的权重体现了本策略的取舍：在碎片化优化场景下，空闲空间是否集中比总空闲比例更关键；但完全忽略总空闲比例也不合理，因为当多个候选都能容纳当前请求时，更高总空闲比例有助于维持segment之间的负载均衡。因此排序采用分层规则，而不是只用单一score：

\begin{enumerate}
  \item \artifact{can\_fit=true}的候选优先于\artifact{can\_fit=false}的候选。
  \item 同为可容纳或同为不可容纳时，\artifact{score}高者优先。
  \item \artifact{score}相近时，\artifact{largest\_free\_region}大者优先。
  \item 仍无法区分时，\artifact{free\_ratio}高者优先。
  \item 最后使用稳定tie-break，避免同一输入下排序非确定。
\end{enumerate}

\begin{lstlisting}[language=C++,caption={\strategy{}候选排序伪代码}]
for candidate in sampled_segments:
    total_free = sum_free_bytes(candidate)
    total_capacity = sum_capacity(candidate)
    largest = largest_contiguous_free_region(candidate)
    free_ratio = total_free / total_capacity
    contiguity = largest / total_free
    can_fit = largest >= request_size
    score = 0.70 * contiguity + 0.30 * free_ratio

sort candidates by:
    can_fit desc,
    score desc,
    largest desc,
    free_ratio desc,
    stable_index asc
\end{lstlisting}

\subsection{权重选择的工程解释}

评分公式中\artifact{contiguity\_ratio}权重为0.70，\artifact{free\_ratio}权重为0.30。这个选择不是为了追求复杂模型，而是为了表达一个工程优先级：在碎片化问题上，连续空间形态应当成为主要信号，但总空闲比例仍然不能被完全丢弃。如果只按最大连续空闲区域排序，策略可能长期偏向少数拥有大连续块的segment，导致整体空间负载不均衡；如果只按总空闲比例排序，则会回到\baseline{}的误判问题。0.70和0.30的组合让策略在两者之间形成明确偏好：先保证当前请求可放入，再兼顾空间均衡。

更重要的是，\strategy{}并不直接用score替代所有排序条件。它先比较\artifact{can\_fit}，再比较score。这意味着当候选A无法容纳当前请求、候选B可以容纳当前请求时，即使A的总空闲比例更高，也不会排在B前面。这个设计避免了“高分但不可用”的情况。score只在同一fit状态内部发挥作用，用于区分多个都可用或多个都不可用的候选。相比只看总空闲比例，它能判断当前请求是否存在足够连续空间；相比只看最大连续空闲区域，它仍保留空间均衡信号；相比只看连续性比例，它不会忽略绝对空闲容量。代价是多一次碎片元数据读取和排序比较，但该成本受候选规模约束。

\subsection{候选排序示例}

为了更直观说明策略行为，可以考虑三个候选segment。Segment A总空闲最高，但最大连续块不足；Segment B总空闲中等，最大连续块刚好满足请求；Segment C最大连续块更大，但总空闲比例较低。若当前请求为10MiB，则A虽然总空闲高，但\artifact{can\_fit=false}，会排在B和C之后；B和C都能容纳请求，再根据综合score和最大连续区域排序。这样排序结果反映的是“先避免明显失败，再在可用候选中做质量选择”。

\begin{table}[htbp]
  \centering
  \caption{候选排序示例}
  \label{tab:rank-example}
  \begin{tabularx}{\textwidth}{|B{1.3cm}|r|r|r|C{2.0cm}|Y|}
    \hline
    \thfirst{segment} & \thcell{总空闲} & \thcell{最大连续空闲} & \thcell{请求大小} & \thcell{\artifact{can\_fit}} & \thcell{排序解释} \\ \hline
    A & 20MiB & 8MiB & 10MiB & false & 总空闲最高但不可容纳，排在可容纳候选之后 \\ \hline
    B & 14MiB & 10MiB & 10MiB & true & 可容纳且空闲较均衡，是优先候选之一 \\ \hline
    C & 12MiB & 12MiB & 10MiB & true & 可容纳且连续块更大，与B按score和tie-break比较 \\ \hline
  \end{tabularx}
\end{table}

\subsection{完整策略流程}

图\ref{fig:strategy-flow}展示了\strategy{}的完整流程。这里需要强调两点。第一，preferred segment阶段不进入评分排序，而是沿用原有优先尝试语义；这是为了保证调用方显式指定的位置仍具有最高优先级。第二，评分排序只发生在采样候选上，且排序后仍逐个调用原有allocator尝试分配；如果候选不足以满足副本数，策略继续进入原有fallback路径。

\begin{figure}[htbp]
  \centering
  \begin{tikzpicture}[
    flowstep/.style={draw=moonblue,very thick,rounded corners=2pt,align=center,minimum width=4.2cm,minimum height=0.82cm,fill=moongray},
    decision/.style={draw=moonorange,very thick,diamond,aspect=2.2,align=center,inner sep=1pt,fill=moonorange!10},
    arrow/.style={-{Latex[length=2.6mm]},thick,draw=moonblue},
    node distance=0.75cm
  ]
    \node[flowstep] (pref) {尝试preferred segments};
    \node[flowstep,below=of pref] (sample) {有界采样候选segments};
    \node[flowstep,below=of sample] (metric) {读取total\_free和largest\_free\_region};
    \node[decision,below=of metric] (fit) {can\_fit?};
    \node[flowstep,below left=0.85cm and 1.1cm of fit] (rankfit) {fit候选优先};
    \node[flowstep,below right=0.85cm and 1.1cm of fit] (rankscore) {按score和tie-break排序};
    \node[flowstep,below=2.1cm of fit] (alloc) {按排序尝试分配};
    \node[flowstep,below=of alloc] (fallback) {未满足副本走原fallback};
    \draw[arrow] (pref) -- (sample);
    \draw[arrow] (sample) -- (metric);
    \draw[arrow] (metric) -- (fit);
    \draw[arrow] (fit) -- node[left,font=\small]{是} (rankfit);
    \draw[arrow] (fit) -- node[right,font=\small]{否} (rankscore);
    \draw[arrow] (rankfit) |- (alloc);
    \draw[arrow] (rankscore) |- (alloc);
    \draw[arrow] (alloc) -- (fallback);
  \end{tikzpicture}
  \caption{\strategy{}策略流程}
  \label{fig:strategy-flow}
\end{figure}

\subsection{多副本场景下的处理}

Mooncake Store并不只处理单副本写入。多副本场景下，策略需要同时满足“每个副本找到可用空间”和“同一次请求尽量不要重复使用同一segment”两个条件。若只为第一个副本选择最优segment，而忽略后续副本，可能导致后续副本候选空间不足。因此，\strategy{}在每次成功分配后都会将该segment加入used集合。后续副本采样和分配时会跳过used segment，从而保留已有策略中的副本去重语义。

这种处理方式也解释了为什么策略没有输出一次性全局排序计划。Store中的空间状态可能在并发请求之间快速变化，提前构造完整计划并不能保证全部成功。更稳妥的方式是：每个副本先读取excluded和used集合，优先尝试preferred segment；随后对未排除、未使用的候选进行碎片感知排序；分配成功后立即把segment加入used集合；如果候选不足，则进入原有fallback路径。这样可以在不引入复杂全局锁和事务计划的情况下，保持较好的工程鲁棒性。

\subsection{退化行为与边界保护}

一个可合入上游的策略必须考虑不完整信息。并不是所有allocator都一定能提供精确的碎片信息；即使能够提供，评分后到实际分配之间也可能发生并发状态变化。因此，\strategy{}从一开始就把碎片信息定位为“排序提示”，而不是“成功承诺”。当\artifact{getLargestFreeRegion()}不可用或返回未知值时，策略使用可用空闲空间进行保守估计；当排序后实际分配失败时，策略不会强行占用空间，而是继续尝试后续候选或fallback。

这种设计有两个好处。第一，它避免新策略对所有allocator形成硬依赖，降低引入成本。第二，它把正确性仍然放在allocator最终检查上，避免因为元数据读取不精确导致空间一致性问题。换句话说，\strategy{}优化的是“尝试顺序”，不是“容量判定权”。容量判定权仍然属于底层allocator。

\subsection{复杂度与并发安全分析}

\strategy{}没有引入新的全局数据结构，也没有在策略内部缓存跨请求状态。它只在每次分配时对采样候选读取allocator元数据，并对候选数组排序。若候选数量为\(k\)，则额外计算复杂度为\(O(k)\)，排序复杂度为\(O(k\log k)\)。由于\(k\)来自现有有界采样，而不是全量segment数\(N\)，策略开销不会随集群规模线性扫描所有segment。

并发方面，新策略不改变allocator的实际申请路径，也不绕过已有锁和状态检查。评分阶段读取的空闲信息只能作为排序依据，不能作为最终成功承诺；真正分配仍由allocator执行。因此，即使评分后其他线程改变了某个segment的空闲状态，最坏结果也只是该候选分配失败并继续尝试后续候选，不会破坏空间一致性。总体上，候选规模沿用有界采样以避免全局扫描，元数据读取只服务于排序，排序成本受候选数量控制；无精确碎片信息时退化为近似总空闲视角，不强制改造所有allocator。

\section{实现路径}

\subsection{代码改动总览}

当前PR补丁涉及9个文件，覆盖策略实现、类型系统、配置解析、启动参数、测试、benchmark、文档和CI辅助步骤。表\ref{tab:files}列出主要改动。

\begin{table}[htbp]
  \centering
  \caption{关键代码改动}
  \label{tab:files}
  \small
  \begin{tabular}{|C{4.2cm}|L{8.5cm}|}
    \hline
    \thfirst{文件} & \thcell{作用} \\ \hline
    \splitpath{mooncake-store/include}{allocation\_strategy.h} & 新增\artifact{FragmentationAwareAllocationStrategy}，实现候选评分、排序、preferred处理和fallback衔接。 \\ \hline
    \splitpath{mooncake-store/include}{types.h} & 新增\artifact{FRAGMENTATION\_AWARE}枚举值，让策略进入统一类型系统。 \\ \hline
    \splitpath{mooncake-store/include}{master\_config.h} & 解析\artifact{allocation\_strategy=fragmentation\_aware}，非法值继续走原有报错路径。 \\ \hline
    \splitpath{mooncake-store/src}{master.cpp} & 更新启动参数帮助信息，确保用户可以从命令行发现新策略。 \\ \hline
    \splitpath{mooncake-store/tests}{allocation\_strategy\_test.cpp} & 增加碎片化定向测试、preferred segment测试和参数化覆盖。 \\ \hline
    \splitpaththree{mooncake-store}{benchmarks}{allocation\_strategy\_bench.cpp} & 将新策略加入分配策略benchmark矩阵，便于后续真实环境对比。 \\ \hline
    \splitpath{docs/source/design}{mooncake-store.md} & 增加策略说明、适用场景和与其他策略的选择建议。 \\ \hline
    \splitpaththree{docs/source/deployment}{mooncake-store-deployment}{-guide.md} & 增加部署参数说明，给出启用方式。 \\ \hline
    \splitpath{.github/workflows}{ci.yml} & 在Go store binding集成测试前释放runner磁盘空间，避免无关环境失败。 \\ \hline
  \end{tabular}
\end{table}

\subsection{策略类接入}

核心实现位于\artifact{mooncake-store/include/allocation\_strategy.h}。新类\artifact{FragmentationAwareAllocationStrategy}继承现有随机策略基础能力，而不是重写所有分配流程。这样做有两个好处：第一，可以复用原有单segment分配、随机fallback和排除集合处理；第二，可以把新增逻辑限制在候选评分和排序层，降低对主路径的侵入。

实现上，新策略先按原有语义处理preferred segment。如果preferred segment没有被排除、没有被当前请求重复使用，并且allocator能够完成申请，则直接返回该分配结果。只有在preferred阶段无法满足全部副本时，策略才进入候选采样和碎片感知排序。排序后的候选按顺序尝试分配，成功后加入used set，避免同一次多副本请求重复选择同一segment。

\subsection{核心执行流程}

从代码路径看，\strategy{}的核心执行可以分为五个阶段。第一阶段是预处理，把excluded segment和当前请求中已经成功使用的segment作为不可选集合。第二阶段是preferred尝试，按调用方传入顺序处理preferred segment。第三阶段是候选采样，沿用已有策略体系的有界采样思想，从剩余segment中得到候选集合。第四阶段是候选评分，对每个候选读取总容量、总空闲、最大连续空闲区域，并计算\artifact{can\_fit}和score。第五阶段是按排序结果尝试实际分配，成功则记录结果，失败则继续尝试后续候选；如果仍不能满足副本数，则回到原有fallback路径。

\begin{lstlisting}[language=C++,caption={\strategy{}实现层执行逻辑概览}]
Allocate(manager, slice_length, replica_num, preferred, excluded):
    used = {}
    result = []

    for name in preferred:
        if name in excluded or name in used:
            continue
        buffer = allocateSingle(manager, name, slice_length)
        if buffer.ok:
            result.push(buffer)
            used.insert(name)
        if result.size == replica_num:
            return result

    candidates = sampleCandidateSegments(manager, excluded, used)
    ranked = []
    for name in candidates:
        metrics = collectFragmentationMetrics(manager, name)
        ranked.push(scoreCandidate(name, metrics, slice_length))

    sort(ranked, FragmentationAwareComparator)

    for candidate in ranked:
        if candidate.name in excluded or candidate.name in used:
            continue
        buffer = allocateSingle(manager, candidate.name, slice_length)
        if buffer.ok:
            result.push(buffer)
            used.insert(candidate.name)
        if result.size == replica_num:
            return result

    return randomBestEffortFallback(manager, result, excluded, used)
\end{lstlisting}

这段逻辑体现了一个重要原则：新策略并没有接管底层空间管理，也没有绕过已有\artifact{allocateSingle}路径。preferred阶段保持调用方显式放置优先级，采样阶段沿用现有有界候选集合，评分阶段把总空闲、容量和最大连续空闲区域转化为排序字段，排序阶段按\artifact{can\_fit}、score和tie-break组织尝试顺序，分配阶段仍调用原有allocator，fallback阶段继续保留best-effort副本语义。这样做能最大限度降低行为风险，因为空间是否真的可分配仍由原有allocator确认；同时，新策略又能在排序阶段规避最明显的碎片化误判。

\subsection{配置和工厂接入}

为了让策略成为正式可选项，而不是测试代码中的孤立类，PR在类型、配置和工厂三处同时接入：

\begin{enumerate}
  \item 在\artifact{AllocationStrategyType}中新增\artifact{FRAGMENTATION\_AWARE}枚举值。
  \item 在\artifact{master\_config.h}中识别字符串\artifact{fragmentation\_aware}。
  \item 在策略工厂中为\artifact{FRAGMENTATION\_AWARE}返回新策略实例。
  \item 在\artifact{master.cpp}和部署文档中更新\artifact{--allocation\_strategy}可选值。
\end{enumerate}

新策略通过Master启动参数显式启用：

\begin{lstlisting}[language=bash]
mooncake_master --allocation_strategy=fragmentation_aware
\end{lstlisting}

由于默认策略保持为\artifact{random}，现有部署不会因为合入该PR而改变行为。若线上灰度时发现新策略不适合某类负载，只需要切回\artifact{random}或\baseline{}即可回滚。

\subsection{兼容性、灰度与回滚}

开源项目中的性能优化如果默认改变行为，评审风险会明显增大。Mooncake Store面向真实推理服务，用户对默认策略的吞吐、稳定性和资源使用已经形成预期。因此，本项目把\strategy{}设计为显式启用策略，而不是替换默认策略。这样既可以让需要碎片优化的用户主动选择，也避免对现有部署产生隐式影响。

灰度使用时，可以先在测试环境或单个Store实例中启用新策略，观察失败分配次数、fallback次数、写入延迟和空间利用率。如果指标改善，再扩大到更多实例；如果发现某类负载下新策略收益不足或元数据读取成本偏高，只需要把启动参数改回\artifact{random}或\baseline{}即可回滚，不需要回滚二进制、修改客户端协议或迁移已有数据。若需要保留空间均衡对照，可继续使用\artifact{--allocation\_strategy=free\_ratio\_first}作为基线。

\subsection{单元测试设计}

测试不是简单地检查“新策略能运行”，而是针对问题定义构造确定性场景。核心测试创建两个segment：\artifact{fragmented}和\artifact{contiguous}。前者总空闲空间更高，但最大连续区域小于请求大小；后者总空闲空间较低，但最大连续区域可以容纳请求。断言目标是：\strategy{}应该选择\artifact{contiguous}，而不是被更高总空闲比例误导。

另一个关键测试是preferred segment语义。即使某个preferred segment在评分上不是最优，只要它没有被排除且可以满足请求，新策略仍应优先选择它。这一点对兼容性很重要，因为preferred segment通常来自上层调度或副本放置约束，不能被内部评分逻辑无条件覆盖。

碎片化测试的构造要避免一个常见陷阱：如果释放的是靠近尾部的块，allocator可能把释放块和尾部剩余空间合并，导致最大连续空闲区域大于预期。这样测试看似构造了碎片，实际却没有形成“大量空闲但连续块不足”的状态。最终PR中的测试选择释放内部块，确保\artifact{fragmented}的总空闲空间较高，但最大连续空闲区域仍小于请求大小；同时构造\artifact{contiguous}，使其总空闲空间较低但最大连续区域满足请求。这种测试更接近问题定义，能稳定验证排序信号是否正确。覆盖点包括碎片化选择、连续空间优先、preferred保留、参数化基础分配语义和benchmark接入。

\subsection{测试修复与CI闭环}

PR创建后经历了三类CI问题并逐步修复。第一次失败来自代码格式检查，后续通过格式化修复。第二次失败来自GitHub runner磁盘不足，Go store binding集成测试在链接阶段出现空间不足；这不是策略逻辑失败，但会阻断PR验证，因此在对应job前增加了磁盘清理步骤。第三次失败来自最初碎片化测试构造不够严格：释放块与尾部空闲区域发生合并，导致最大连续区域大于预期，断言不稳定。后续测试改为构造5个8MiB块并释放内部块，避免尾部合并影响碎片形态。上述问题修复后重新触发CI，PR头提交\artifact{0123fa1}通过26个检查，1个检查跳过。

\subsection{文档、benchmark与复现材料}

除了代码本身，PR和比赛材料还补充了三类工程化内容。第一，在Mooncake官方文档中说明\strategy{}的适用场景和启用方式，避免评审只看到代码而不知道何时使用。第二，将新策略加入\artifact{allocation\_strategy\_bench.cpp}，让后续真实Store benchmark可以直接横向比较\artifact{random}、\baseline{}和\strategy{}。第三，在比赛材料仓库中提供独立复现实验和日志，用轻量仿真固定问题模型，证明改进来自“排序信号变化”，而不是扩大候选规模或修改实验口径。

\section{验证与效果}

\subsection{验证思路}

本项目的验证分为三层。第一层是代码级验证，确认新策略能够编译、能够被配置解析、能够进入策略工厂，并且不会破坏已有分配策略测试。第二层是问题级验证，通过最小碎片化测试证明策略确实解决“总空闲比例高但最大连续空闲不足”的误判。第三层是专题仿真验证，通过多个请求和多个segment构造，观察首选成功次数、fallback尝试次数和单次决策成本。三层验证分别回答“能不能运行”“是否解决目标问题”“收益和成本如何”。

需要说明的是，初赛阶段的验证重点是策略层正确性和可复现性，而不是宣称已经完成完整生产压测。真实RDMA环境、真实SGLang HiCache端到端吞吐和长时间集群稳定性测试需要更多硬件和部署条件，适合作为决赛阶段继续推进。本报告中的数据口径均围绕策略层展开，不把离线仿真结果夸大为生产集群吞吐结论。代码级验证依赖单元测试、参数化测试和上游CI；问题级验证依赖碎片化定向测试；指标级验证依赖专题仿真和日志；工程级验证依赖文档、benchmark接入和PR流程。

\subsection{CI验证}

Mooncake上游draft PR状态如表\ref{tab:ci}所示。

\begin{table}[h]
  \centering
  \caption{PR CI状态}
  \label{tab:ci}
  \begin{tabularx}{\textwidth}{|B{3.2cm}|Y|}
    \hline
    \thfirst{项目} & \thcell{内容} \\ \hline
    PR链接 & \url{https://github.com/kvcache-ai/Mooncake/pull/2797} \\ \hline
    头提交 & \artifact{0123fa1 Fix fragmentation-aware allocation test setup} \\ \hline
    CI结果 & 26个检查成功，1个检查跳过 \\ \hline
    结论 & 上游构建、格式检查、文档构建和相关测试已通过。 \\ \hline
  \end{tabularx}
\end{table}

\subsection{专题仿真指标}

专题仿真使用5个segment和6个大对象请求，模拟高总空闲但碎片化、低总空闲但连续空间充足两类segment并存的Store状态。结果如表\ref{tab:metrics}所示。

这里的“首选segment可直接容纳”指排序后的第一个候选是否能够立即满足当前请求。这个指标非常关键，因为它衡量的是策略排序质量，而不是系统最终是否能通过fallback补救。“最终可容纳”表示在继续尝试其他候选后，请求最终是否能找到空间；它用于确认两个策略面对同一容量条件时并没有改变总可用容量。“fallback尝试次数”则用于衡量首选错误带来的额外路径压力。平均候选segment数量用于证明新策略没有通过扩大搜索范围获得收益，而是在相同候选规模内改进排序。

\begin{table}[h]
  \centering
  \caption{专题对齐仿真结果}
  \label{tab:metrics}
  \begin{tabular}{|C{5.6cm}|r|r|}
    \hline
    \thfirst{指标} & \thcell{\baseline{}} & \thcell{\strategy{}} \\ \hline
    大对象首选segment可直接容纳 & 0/6 & 6/6 \\ \hline
    最终可容纳 & 6/6 & 6/6 \\ \hline
    fallback尝试次数 & 11 & 0 \\ \hline
    平均候选segment数量 & 5.00 & 5.00 \\ \hline
    单次决策耗时 & 121.00ns & 138.93ns \\ \hline
    新增决策开销 & 不适用 & 17.93ns \\ \hline
  \end{tabular}
\end{table}

\begin{figure}[h]
  \centering
  \resizebox{0.95\textwidth}{!}{%
  \begin{tikzpicture}[
    axis/.style={draw=black!60},
    barA/.style={fill=moonred!78},
    barB/.style={fill=moongreen!78},
    label/.style={font=\small}
  ]
    \draw[axis] (0,0) -- (4.4,0);
    \draw[axis] (0,0) -- (0,3.6);
    \foreach \y/\txt in {0/0,1.65/3,3.3/6} {
      \draw[black!30] (0,\y) -- (4.4,\y);
      \node[label,anchor=east] at (-0.15,\y) {\txt};
    }
    \path[barA] (1.0,0) rectangle (1.9,0.02);
    \path[barB] (2.5,0) rectangle (3.4,3.3);
    \node[label] at (1.45,0.28) {0/6};
    \node[label] at (2.95,3.55) {6/6};
    \node[label,align=center] at (1.45,-0.35) {基线};
    \node[label,align=center] at (2.95,-0.35) {新策略};
    \node[font=\small\bfseries] at (2.2,-0.85) {首选成功次数};

    \draw[axis] (5.7,0) -- (10.1,0);
    \draw[axis] (5.7,0) -- (5.7,3.6);
    \foreach \y/\txt in {0/0,1.65/6,3.3/12} {
      \draw[black!30] (5.7,\y) -- (10.1,\y);
      \node[label,anchor=east] at (5.55,\y) {\txt};
    }
    \path[barA] (6.7,0) rectangle (7.6,3.03);
    \path[barB] (8.2,0) rectangle (9.1,0.02);
    \node[label] at (7.15,3.28) {11};
    \node[label] at (8.65,0.28) {0};
    \node[label,align=center] at (7.15,-0.35) {基线};
    \node[label,align=center] at (8.65,-0.35) {新策略};
    \node[font=\small\bfseries] at (7.9,-0.85) {fallback尝试次数};

    \node[label,align=center] at (5.05,-1.35)
      {基线 = \baseline{}；新策略 = \strategy{}};
  \end{tikzpicture}
  }
  \caption{首选成功次数与fallback尝试次数对比}
\end{figure}

\subsection{结果解释}

该结果有三个含义。第一，\strategy{}显著改善首选segment质量，在所有测试请求上都能首先选中可容纳segment。第二，平均候选segment数量没有增加，说明改进来自排序信号变化，而不是扩大搜索范围。第三，单次决策新增开销约17.93ns，在本地CPU仿真中处于很小量级，符合“用少量元数据计算换取更稳定分配路径”的设计目标。

进一步看，\baseline{}最终也能达到6/6可容纳，说明测试场景中系统总容量并不缺失；问题在于它先尝试了不合适的segment，需要额外fallback。换言之，本项目优化的不是“凭空增加容量”，而是“更快找到已经存在的正确位置”。这也是Store分配策略优化的现实价值：在容量足够但碎片化严重的状态下，吞吐和尾延迟问题往往来自无效尝试、重试和锁路径放大，而不一定来自物理资源绝对不足。

单次决策耗时从121.00ns增加到138.93ns，新增约17.93ns。这个成本主要来自读取最大连续空闲区域、计算连续性比例和增加排序比较字段。与一次真实网络I/O、远程存储访问或复杂推理请求相比，这一级别开销很小；但报告仍然保留该成本，而不把策略描述为“零开销”。更准确的结论是：\strategy{}用可控的纳秒级决策开销，换取碎片化场景下明显更高的首选成功率和更少fallback尝试。

\section{与赛题要求的对应}

\begin{longtable}{|B{3.2cm}|L{5.0cm}|L{5.0cm}|}
  \caption{赛题要求映射}\\
  \hline
\thfirst{赛题关注点} & \thcell{本项目实现} & \thcell{当前证据} \\ \hline
  Mooncake Store吞吐性能 & 减少碎片化导致的无效首选分配和fallback压力 & 专题仿真中fallback尝试次数从11降为0 \\ \hline
  可扩展性 & 保持有界候选采样，不做全局扫描 & 平均候选segment数量保持5.00 \\ \hline
  高可用语义 & 不改变副本best-effort路径，保留preferred/excluded语义 & 单元测试和策略设计说明 \\ \hline
  SGLang HiCache+Store & 优化Store后端对象放置路径，为HiCache后端稳定性提供基础 & 方案定位和代码路径位于Mooncake Store \\ \hline
  开源规范 & 以PR形式提交上游Mooncake，并保留patch和文档 & PR\#2797通过CI \\ \hline
  可复现性 & 提供仿真源码、日志、技术报告和使用说明 & \artifact{repro/}、\artifact{logs/}、\artifact{EVALUATION.md} \\ \hline
\end{longtable}

\section{使用与复现}

\subsection{启用策略}

\begin{lstlisting}[language=bash]
mooncake_master --allocation_strategy=fragmentation_aware
\end{lstlisting}

\subsection{应用补丁}

\begin{lstlisting}[language=bash]
git apply --check mooncake_fragmentation_aware_pr_2797_0123fa1.patch
git apply mooncake_fragmentation_aware_pr_2797_0123fa1.patch
\end{lstlisting}

\subsection{回滚补丁}

\begin{lstlisting}[language=bash]
git apply -R mooncake_fragmentation_aware_pr_2797_0123fa1.patch
\end{lstlisting}

\subsection{复现实验}

\begin{lstlisting}[language=bash]
g++ -std=c++17 -O2 repro/fragmentation_aware_sim.cpp -o fragmentation_aware_sim
./fragmentation_aware_sim

g++ -std=c++17 -O2 repro/fragmentation_aware_metrics.cpp -o fragmentation_aware_metrics
./fragmentation_aware_metrics

g++ -std=c++17 -O2 repro/topic_aligned_store_scalability_sim.cpp -o topic_aligned_store_scalability_sim
./topic_aligned_store_scalability_sim
\end{lstlisting}

\section{边界与后续工作}

本项目当前阶段已经完成代码实现、上游PR、CI验证、仿真评估和中文提交材料整理，但仍有边界需要明确：

\begin{itemize}
  \item 当前不宣称已经完成真实RDMA硬件benchmark。
  \item 当前不宣称已经完成真实SGLang HiCache端到端吞吐压测。
  \item 当前不宣称已经重设计Mooncake高可用协议。
  \item 当前专题仿真验证的是分配策略层，不等同于完整生产集群压测。
\end{itemize}

后续如果进入决赛，建议补充三类工作：第一，在真实Mooncake Store环境中运行\artifact{allocation\_strategy\_bench}，统计吞吐、P50/P99、失败分配次数、fallback次数和内存利用率；第二，在具备RDMA条件的集群中验证跨节点场景；第三，将SGLang HiCache接入Mooncake Store后端，使用真实推理负载验证KVCache复用和后端IO表现。

\section{结论}

本文围绕Mooncake赛题2提出并实现Mooncake Store碎片感知分配策略。该策略针对\baseline{}只看总空闲比例的不足，引入最大连续空闲区域和可容纳性排序，在不改变默认行为、不扩大候选规模、不破坏原有副本语义的前提下，提升碎片化场景下的大对象首选分配质量。当前PR已经通过上游GitHub Actions，确定性仿真显示首选成功次数从0/6提升到6/6，fallback尝试次数从11降到0。该工作是一个边界清晰、可复现、可回滚、适合继续扩展到真实Store benchmark和SGLang HiCache端到端测试的初赛成果。

\appendix
\section{提交材料清单}

\begin{longtable}{|C{6.2cm}|L{7.2cm}|}
  \hline
\thfirst{文件或目录} & \thcell{说明} \\ \hline
  \artifact{README.md} & 项目入口说明 \\ \hline
  \artifact{SUBMISSION.md} & 平台提交说明 \\ \hline
  \artifact{DESIGN.md} & 设计文档 \\ \hline
  \artifact{technical\_solution.md} & 技术方案说明 \\ \hline
  \artifact{EVALUATION.md} & 评估报告 \\ \hline
  \artifact{testing.md} & 测试说明 \\ \hline
  \artifact{usage.md} & 使用说明 \\ \hline
  \artifact{REVIEW\_GUIDE.md} & 评审阅读指南 \\ \hline
  \splitpath{mooncake\_fragmentation\_aware}{pr\_2797\_0123fa1.patch} & 当前PR补丁 \\ \hline
  \artifact{repro/} & 复现实验源码 \\ \hline
  \artifact{logs/} & 验证日志 \\ \hline
  \artifact{slides/} & 初赛展示PPT \\ \hline
  \artifact{report/} & 本技术报告源码和PDF \\ \hline
  \artifact{release/} & 提交压缩包和校验文件 \\ \hline
\end{longtable}

\section{参考链接}

\begin{itemize}
  \item Mooncake上游仓库：\href{https://github.com/kvcache-ai/Mooncake}{kvcache-ai/Mooncake}
  \item 本项目Mooncake PR：\href{https://github.com/kvcache-ai/Mooncake/pull/2797}{kvcache-ai/Mooncake\#2797}
  \item 比赛材料仓库：\href{https://github.com/Lorry1024/ccf2026-mooncake-fragmentation-aware}{Lorry1024/ccf2026-mooncake-fragmentation-aware}
  \item Mooncake fork分支：\href{https://github.com/Lorry1024/Mooncake/tree/ccf-fragmentation-aware-allocation}{Lorry1024/Mooncake:ccf-fragmentation-aware-allocation}
\end{itemize}

\end{document}
